%%
%% This is file `coppe-quickref.tex',
%% generated with the docstrip utility.
%%
%% The original source files were:
%%
%% coppe.dtx  (with options: `quickref')
%% 
%% This is one of the demonstration documents of the `coppe' document class:
%% the five main-language demos, the PDF/A build of the full example and the
%% side-by-side cover sheet.  It is generated from coppe.dtx by docstrip; edit
%% the .dtx source, not this file.
%% 
%% \CheckSum{0}
%% \CharacterTable
%%  {Upper-case    \A\B\C\D\E\F\G\H\I\J\K\L\M\N\O\P\Q\R\S\T\U\V\W\X\Y\Z
%%   Lower-case    \a\b\c\d\e\f\g\h\i\j\k\l\m\n\o\p\q\r\s\t\u\v\w\x\y\z
%%   Digits        \0\1\2\3\4\5\6\7\8\9
%%   Exclamation   \!     Double quote  \"     Hash (number) \#
%%   Dollar        \$     Percent       \%     Ampersand     \&
%%   Acute accent  \'     Left paren    \(     Right paren   \)
%%   Asterisk      \*     Plus          \+     Comma         \,
%%   Minus         \-     Point         \.     Solidus       \/
%%   Colon         \:     Semicolon     \;     Less than     \<
%%   Equals        \=     Greater than  \>     Question mark \?
%%   Commercial at \@     Left bracket  \[     Backslash     \\
%%   Right bracket \]     Circumflex    \^     Underscore    \_
%%   Grave accent  \`     Left brace    \{     Vertical bar  \|
%%   Right brace   \}     Tilde         \~}
%%
%% coppe-quickref.tex -- CoppeTeX quick reference, in English.
%%
%% GERADO de coppe.dtx. Nao edite este arquivo: edite o modulo `quickref' do
%% coppe.dtx e rode `pdflatex coppe.ins'.
\documentclass[a4paper,10pt]{article}
\usepackage[T1]{fontenc}
\usepackage[utf8]{inputenc}
\usepackage{lmodern}
\usepackage[margin=2cm,landscape]{geometry}
\usepackage{longtable}
\usepackage{booktabs}
\usepackage{array}
\usepackage{xcolor}
\usepackage[colorlinks=true,allbordercolors=white]{hyperref}

\newcommand\coppeversion{v4.1}

\newcolumntype{C}{>{\ttfamily\raggedright\arraybackslash}p{0.20\textwidth}}
\newcolumntype{E}{>{\ttfamily\footnotesize\raggedright\arraybackslash}p{0.30\textwidth}}
\newcolumntype{W}{>{\centering\arraybackslash}p{0.055\textwidth}}
\newcolumntype{D}{>{\raggedright\arraybackslash}p{0.37\textwidth}}

\newcommand\pre{\textsc{pre}}
\newcommand\doc{\textsc{doc}}

\setlength\parindent{0pt}

\begin{document}

\begin{center}
  {\Large\bfseries CoppeTeX quick reference}\\[2pt]
  {\normalsize Class \texttt{coppe}, \coppeversion\ --- every option, command and
  environment}\\[6pt]
  \begin{minipage}{0.86\textwidth}\small
  This is the short companion to \texttt{coppe.pdf}, the full manual, which is
  written in Portuguese. A COPPE thesis may be written in Portuguese, English or
  Spanish (CEPG Resolution 302/2024, art.~57); this sheet is for whoever writes
  it in English.

  \medskip
  The \textbf{Where} column says where a command may appear:
  \pre\ means the preamble, before \verb|\begin{document}|;
  \doc\ means after \verb|\begin{document}|, and, for everything that feeds the
  cover, \emph{before} \verb|\maketitle|.

  \medskip
  \textbf{Three things that surprise people.} The document data goes
  \emph{inside} \texttt{document}, not in the preamble. The bibliography is
  \texttt{biblatex} with \texttt{biber}, not Bib\TeX. And
  \verb|\frontmatter|/\verb|\mainmatter| are what switch folio numbering on, so
  leaving them out gives the wrong page numbers.
  \end{minipage}
\end{center}

\vspace{4pt}

\begin{longtable}{@{}CEWD@{}}
\toprule
\normalfont\bfseries Command & \normalfont\bfseries Example &
\normalfont\bfseries Where & \normalfont\bfseries What it does \\
\midrule
\endfirsthead
\toprule
\normalfont\bfseries Command & \normalfont\bfseries Example &
\normalfont\bfseries Where & \normalfont\bfseries What it does \\
\midrule
\endhead
\bottomrule
\endfoot

\multicolumn{4}{@{}l@{}}{\normalfont\bfseries\small CLASS OPTIONS --- work
  type, one is required}\\[2pt]
msc      & \textbackslash documentclass[msc]\{coppe\}     & \pre & M.Sc.\ dissertation. \\
dsc      & \textbackslash documentclass[dsc]\{coppe\}     & \pre & D.Sc.\ thesis. \\
mscexam  & \textbackslash documentclass[mscexam]\{coppe\} & \pre & M.Sc.\ qualifying exam. No CAPES sheet, no catalogue card, no reference above the abstract. \\
dscexam  & \textbackslash documentclass[dscexam]\{coppe\} & \pre & D.Sc.\ qualifying exam. Same. \\
mscsem   & \textbackslash documentclass[mscsem]\{coppe\}  & \pre & M.Sc.\ seminar, a requirement of some Programas and \emph{not} a qualifying exam. Same three omissions: it is not deposited. \\
\midrule

\multicolumn{4}{@{}l@{}}{\normalfont\bfseries\small CLASS OPTIONS --- writing
  language}\\[2pt]
brazilian & [brazilian,dsc] & \pre & Portuguese. The default; needs no writing. \\
english   & [english,dsc]   & \pre & English. The foreign abstract switches to Portuguese. \\
spanish   & [spanish,dsc]   & \pre & Spanish. Needs \texttt{coppe-lang-spanish.def} and \texttt{spanish-coppe.lbx}. Three abstracts. \\
french    & [french,dsc]    & \pre & French. \textbf{Not} a permitted thesis language; ships as a demonstration. \\
italian   & [italian,dsc]   & \pre & Italian. Same. \\
\midrule

\multicolumn{4}{@{}l@{}}{\normalfont\bfseries\small CLASS OPTIONS --- the
  rest}\\[2pt]
numbers   & [dsc,numbers]   & \pre & Numeric citations \texttt{[1]} instead of author-date; list in citation order. \\
pdfa      & [dsc,pdfa]      & \pre & PDF/A-2b output, required for deposit (2.2d). Turn it on when you deposit. \\
comserifa & [dsc,comserifa] & \pre & Serif family. Sans is the default since v4.1. \\
semserifa & ---             & \pre & Accepted, does nothing. What it asked for became the default. \\
semlinks  & [dsc,semlinks]  & \pre & Links lose colour and frame; they stay clickable. For printing. \\
doublespacing & [dsc,doublespacing] & \pre & Double spacing instead of one-and-a-half (2.2c). \\
twoside   & [dsc,twoside]   & \pre & Mirrored margins, for printing and binding. \\
oneside   & [dsc,oneside]   & \pre & One-sided, 3\,cm left margin on every sheet. The default. \\
orientadorexamina & [dsc,orientadorexamina] & \pre & Put advisors back on the approval sheet. Off by default: the sheet records who examined. \\
coorientador & [dsc,coorientador] & \pre & List coadvisors on the abstract pages too. \\
rascunhoficha & [dsc,rascunhoficha] & \pre & Draft catalogue card while writing. \textbf{Never valid for deposit.} \\
resumosemreferencia & [dsc,resumosemreferencia] & \pre & Drop the work's own reference from the top of the three abstracts. It is printed by default (3.1.2.1.4, Annexes E and F). \\
listasnosumario & [dsc,listasnosumario] & \pre & Put the pre-textual lists back in the table of contents. \\
assinaturas & ---           & \pre & Accepted, does nothing. Signature rules are gone: deposit is digital. \\
numeraisromanos & [dsc,numeraisromanos] & \pre & Roman numerals on pre-textual sheets. \textbf{Against 2.7}; legacy only. \\
morewrites & [dsc,morewrites] & \pre & A missing \texttt{morewrites} package becomes an error, not a warning. \\
semmorewrites & [dsc,semmorewrites] & \pre & Do not load \texttt{morewrites} at all. \\
\midrule

\multicolumn{4}{@{}l@{}}{\normalfont\bfseries\small PREAMBLE}\\[2pt]
\textbackslash addbibresource & \textbackslash addbibresource\{thesis.bib\} & \pre & Declare the bibliography database. With the \texttt{.bib} extension. \\
\textbackslash makelosymbols & \textbackslash makelosymbols & \pre & Turn on collection of the list of symbols. Without it, \texttt{\textbackslash symbl} records nothing. \\
\textbackslash makeloabbreviations & \textbackslash makeloabbreviations & \pre & Same for the list of abbreviations. \\
\textbackslash newcoppefloat & \textbackslash newcoppefloat\{map\}\{Map\}\{List of Maps\} & \pre & Declare your own float: caption on top, \texttt{\textbackslash source} below, own list. \\
\textbackslash newsigla & \textbackslash newsigla\{ufrj\}\{UFRJ\}\{Universidade Federal...\} & \pre & Declare an acronym for \texttt{\textbackslash sigla}. \\
\textbackslash usecoppelanguage & \textbackslash usecoppelanguage\{french\} & \pre & Load one more language pack, for quotations in a third language. \\
\textbackslash graphicspath & \textbackslash graphicspath\{\{figures/\}\} & \pre & Where your figures live. Standard \LaTeX. \\
\textbackslash freeconfig & \textbackslash freeconfig\{Dr.\}\{...\}\{...\}\{...\}\{...\}\{...\} & \pre & Override the work type entirely. Escape hatch; do not use for a normal thesis. \\
\midrule

\multicolumn{4}{@{}l@{}}{\normalfont\bfseries\small IDENTIFICATION --- after
  \textbackslash begin\{document\}, before \textbackslash maketitle}\\[2pt]
\textbackslash title & \textbackslash title\{Título da tese\} & \doc & The Portuguese title. \\
\textbackslash foreigntitle & \textbackslash foreigntitle\{Thesis title\} & \doc & The English title. \\
\textbackslash titlein & \textbackslash titlein\{spanish\}\{Título\} & \doc & The title in any other language. \\
\textbackslash subtitle & \textbackslash subtitle\{Um subtítulo\} & \doc & Optional subtitle, Portuguese. Printed after a colon, in ordinary case. \\
\textbackslash foreignsubtitle & \textbackslash foreignsubtitle\{A subtitle\} & \doc & The same, English. \\
\textbackslash subtitlein & \textbackslash subtitlein\{spanish\}\{...\} & \doc & The same, any language. \\
\textbackslash author & \textbackslash author\{All given names\}\{Surname\} & \doc & Two arguments: the catalogue card has to invert the name. \\
\textbackslash advisor & \textbackslash advisor\{Ana\}\{Lima\}\{D.Sc.\}\{UFRJ\} & \doc & Name, surname, degree, institution. Repeat per advisor. Optional first argument is the treatment. \\
\textbackslash coadvisor & \textbackslash coadvisor\{Bia\}\{Sousa\}\{Ph.D.\}\{UFF\} & \doc & Same four arguments. \\
\textbackslash examiner & \textbackslash examiner\{Carlos Dias\}\{D.Sc.\}\{USP\} & \doc & Name and surname in one argument, then degree and institution. \\
\textbackslash department & \textbackslash department\{PESC\} & \doc & The Programa's acronym. The class knows all thirteen. \\
\textbackslash date & \textbackslash date\{09\}\{2026\} & \doc & Month and year of deposit, as numbers. Not the defence date. \\
\textbackslash dataaprovacao & \textbackslash dataaprovacao\{15 de setembro de 2026\} & \doc & Defence date, in words. Without it the sheet reads ``a ser determinada''. \\
\textbackslash keyword & \textbackslash keyword\{Formal methods\} & \doc & One call per keyword, main language. \\
\textbackslash foreignkeyword & \textbackslash foreignkeyword\{Métodos formais\} & \doc & One call per keyword, foreign language. \\
\textbackslash braziliankeyword & \textbackslash braziliankeyword\{Métodos formais\} & \doc & Portuguese keywords when the work is not in Portuguese. \\
\textbackslash volumes & \textbackslash volumes\{3\} & \doc & How many volumes. Nothing is printed when there is one. \\
\textbackslash volume & \textbackslash volume\{1\} & \doc & Which volume this is. \\
\midrule

\multicolumn{4}{@{}l@{}}{\normalfont\bfseries\small TITLE PAGE AND CAPES
  SHEET}\\[2pt]
\textbackslash areaconcentracao & \textbackslash areaconcentracao\{Engenharia de Sistemas\} & \doc & Area of concentration. On three sheets. \\
\textbackslash linhapesquisa & \textbackslash linhapesquisa\{Engenharia de Dados\} & \doc & Research line, on the title page. Label follows the main language. Empty is fine. \\
\textbackslash tipoproducao & \textbackslash tipoproducao\{bibliografica\} & \doc & One of \texttt{bibliografica}, \texttt{artistica}, \texttt{tecnologica}, \texttt{tecnica}. \\
\textbackslash projetovinculado & \textbackslash projetovinculado\{sim\} & \doc & \texttt{sim} or \texttt{nao}: is there a linked research project? \\
\textbackslash nomeprojeto & \textbackslash nomeprojeto\{Nome do projeto\} & \doc & Its name, if there is one. \\
\textbackslash agenciafomento & \textbackslash agenciafomento\{Conselho...\}\{CNPq\} & \doc & Funding agency, full name and acronym. One call each. \\
\textbackslash fichacatalografica & \textbackslash fichacatalografica\{ficha.pdf\} & \doc & The official catalogue card, as a PDF from the SiBI generator. \\
\midrule

\multicolumn{4}{@{}l@{}}{\normalfont\bfseries\small STRUCTURE}\\[2pt]
\textbackslash maketitle & \textbackslash maketitle & \doc & Typesets cover, title page and CAPES sheet. After all the data. \\
\textbackslash frontmatter & \textbackslash frontmatter & \doc & Opens the pre-textual part and typesets the approval sheet. \\
\textbackslash mainmatter & \textbackslash mainmatter & \doc & Opens the textual part. This is what starts printing the folio. \\
\textbackslash backmatter & \textbackslash backmatter & \doc & Opens the post-textual part. \\
\textbackslash dedication & \textbackslash dedication\{To those who waited.\} & \doc & Dedication, on a sheet of its own. \\
abstract & \textbackslash begin\{abstract\} ... & \doc & Abstract in the main language. One sheet, one paragraph, 250 words. \\
foreignabstract & \textbackslash begin\{foreignabstract\} ... & \doc & Abstract in the foreign language. \\
brazilianabstract & \textbackslash begin\{brazilianabstract\} ... & \doc & Third abstract, Portuguese, for a Spanish- or French-main work. \\
\textbackslash appendix & \textbackslash appendix & \doc & Opens the appendices: text \emph{you} wrote. \\
\textbackslash annex & \textbackslash annex & \doc & Opens the annexes: documents you did \emph{not} write. \\
\textbackslash printbibliography & \textbackslash printbibliography & \doc & The reference list, titled ``Referências'' (3.1.4.1). \\
\textbackslash coppetexfinalpage & \textbackslash coppetexfinalpage & \doc & Optional colophon: class version, engine, font. \\
\midrule

\multicolumn{4}{@{}l@{}}{\normalfont\bfseries\small PRE-TEXTUAL LISTS ---
  lists first, table of contents last (3.1.2.1.6)}\\[2pt]
\textbackslash listoffigures & \textbackslash listoffigures & \doc & List of figures. \\
\textbackslash listoftables & \textbackslash listoftables & \doc & List of tables. \\
\textbackslash listofquadros & \textbackslash listofquadros & \doc & List of \emph{quadros} (framed on all sides, unlike a table). \\
\textbackslash listofprogramas & \textbackslash listofprogramas & \doc & List of code listings. \\
\textbackslash listofalgorithms & \textbackslash listofalgorithms & \doc & List of algorithms. \\
\textbackslash printloabbreviations & \textbackslash printloabbreviations & \doc & List of abbreviations and acronyms. No page numbers (4.1.1). \\
\textbackslash printlosymbols & \textbackslash printlosymbols & \doc & List of symbols. Likewise. \\
\textbackslash tableofcontents & \textbackslash tableofcontents & \doc & The table of contents. Comes last of all the lists. \\
\midrule

\multicolumn{4}{@{}l@{}}{\normalfont\bfseries\small ILLUSTRATIONS ---
  caption on top, source below (2.10)}\\[2pt]
\textbackslash source & \textbackslash source\{Prepared by the author.\} & \doc & The source line, below the float. Mandatory by the norm. \\
\textbackslash fonte & \textbackslash fonte\{Elaboração própria.\} & \doc & The same command, Portuguese name. \\
\textbackslash cpsource & \textbackslash cpsource\{...\} & \doc & Clash-safe alias, always available. \\
\textbackslash cpfonte & \textbackslash cpfonte\{...\} & \doc & Likewise. \\
quadro & \textbackslash begin\{quadro\}[ht] ... & \doc & A float framed on all sides. Behaves like \texttt{figure}. \\
python & \textbackslash begin\{python\}[caption=\{...\}] & \doc & Python listing, syntax-highlighted, caption on top. \\
java, xml, html, prolog & \textbackslash begin\{java\}[...] ... & \doc & The same, for those languages. \\
\textbackslash pythoninline & \textbackslash pythoninline\{f(x)\} & \doc & Code inside a paragraph. \\
\textbackslash javainline & \textbackslash javainline\{new X()\} & \doc & Likewise. \\
\textbackslash pythonexternal & \textbackslash pythonexternal\{code.py\} & \doc & Listing read from a file. \\
\textbackslash javaexternal & \textbackslash javaexternal\{Code.java\} & \doc & Likewise. \\
\textbackslash pythonstyle & \textbackslash pythonstyle & \doc & Apply a language's style without opening the environment. Also \texttt{\textbackslash javastyle}, \texttt{\textbackslash xmlstyle}, \texttt{\textbackslash htmlstyle}, \texttt{\textbackslash prologstyle}. \\
\midrule

\multicolumn{4}{@{}l@{}}{\normalfont\bfseries\small QUOTATIONS, ACRONYMS,
  SYMBOLS}\\[2pt]
\textbackslash citacao & \textbackslash citacao\{up to three lines\} & \doc & Short quotation, inside the paragraph, in quotation marks. \\
longquote & \textbackslash begin\{longquote\} ... & \doc & Quotation over three lines: indented 4\,cm past the margin, smaller, single spaced, no quotation marks (2.2b). \\
\textbackslash abbrev & \textbackslash abbrev[IOT]\{IoT\}\{Internet of Things\} & \doc & Register an acronym. Optional first argument is the sort key. \\
\textbackslash symbl & \textbackslash symbl\{\$\textbackslash alpha\$\}\{Learning rate\} & \doc & Register a symbol. Same optional sort key. \\
\textbackslash sigla & \textbackslash sigla\{ufrj\} & \doc & Full form on first use, acronym afterwards (2.8). \texttt{\textbackslash sigla*} forces the full form. \\
\midrule

\multicolumn{4}{@{}l@{}}{\normalfont\bfseries\small CITATIONS --- biblatex,
  with the natbib layer on}\\[2pt]
\textbackslash cite & \textbackslash cite\{key\} & \doc & (SURNAME, 2026), or [1] with \texttt{numbers}. \\
\textbackslash citet & \textbackslash citet\{key\} & \doc & SURNAME (2026). \\
\textbackslash citep & \textbackslash citep\{key\} & \doc & (SURNAME, 2026). \\
\textbackslash textcite & \textbackslash textcite\{key\} & \doc & Same as \texttt{\textbackslash citet}, biblatex spelling. \\
\textbackslash parencite & \textbackslash parencite\{key\} & \doc & Same as \texttt{\textbackslash citep}. \\
\textbackslash citeauthor & \textbackslash citeauthor\{key\} & \doc & The author alone. \\
\textbackslash citeyear & \textbackslash citeyear\{key\} & \doc & The year alone. \\
\textbackslash citetitle & \textbackslash citetitle\{key\} & \doc & The title alone. \\
\textbackslash cite[p.~45] & \textbackslash cite[p.~45, our translation]\{key\} & \doc & Page, and any note, in the optional argument. \\
\textbackslash citepapud & \textbackslash citepapud[45]\{Original\}\{1953\}\{key-read\} & \doc & Citation of a citation (4.1.2.2.1). \textbf{Three} mandatory arguments: the original author and year, which are printed as text, and the key of the work you actually read, which is the only one that enters the list. Also \texttt{\textbackslash citetapud}. \\
\textbackslash nocite & \textbackslash nocite\{key\} & \doc & Into the list without citing in the text. \\

\end{longtable}

\vfill
\begin{center}\small
Full manual: \texttt{coppe.pdf}, in Portuguese.\quad
Complete worked example: \texttt{example.tex}.\quad
Problems: \url{https://github.com/COPPE-UFRJ/CoppeTeX/issues}
\end{center}

\end{document}

%% 
%%
%% End of file `coppe-quickref.tex'.
